% Project UML use-case diagram. Each use-case ID links to its table row.
\begin{tikzpicture}[
  external actor/.style={draw=fsaGray,rounded corners=2pt,fill=fsaGray!8,
    align=center,text width=2.5cm,minimum height=1.15cm,font=\small},
  project use case/.style={ellipse,draw=fsaBlue,fill=fsaLightBlue,
    align=center,text width=3.8cm,minimum height=1.25cm,font=\small},
  association/.style={draw=fsaGray,line width=0.7pt}
]
  % Reserve a separate header band above the use-case ovals.
  \draw[fsaBlue,rounded corners=3pt] (-2.95,-2.0) rectangle (2.95,3.2);
  \node[align=center,font=\small\bfseries,text width=4.8cm] at (0,2.45)
    {Sensor-interface subsystem\\FTDI245R acquisition\\pipeline code base};
  \node[project use case] (acquire) at (0,0.65)
    {\UseCaseRef{UC-001}\\Acquire sensor data};
  \node[project use case] (outcomes) at (0,-1.0)
    {\UseCaseRef{UC-002}\\Handle acquisition\\read outcomes};
  \node[external actor] (engineer) at (-4.5,-0.15)
    {\textless\textless actor\textgreater\textgreater\\Embedded Engineer};
  \node[external actor] (consumer) at (4.6,0.65)
    {\textless\textless actor\textgreater\textgreater\\Downstream processing\\consumer};
  \draw[association] (engineer.east) -- (-3.05,-0.15) |- (acquire.west);
  \draw[association] (engineer.east) -- (-3.05,-0.15) |- (outcomes.west);
  \draw[association] (acquire.east) -- (consumer.west);
\end{tikzpicture}
